\Askhsh{22151_SOLUTION}{1}
\begin{ylh}{1}
    \item [ΛΥΣΗ]
\end{ylh}
\begin{alist}
    \item 
    \begin{rlist}
        \item Το εμβαδόν $E_{\text{ΑΕ}}$ του κύκλου $(A, r)$ είναι ίσο με $E_{\text{ΑΕ}} = \pi \cdot r^2$ και το εμβαδόν $E_{\text{ΕΓ}}$ του δακτυλίου μεταξύ των κύκλων $(A, r)$ και $(A, R)$ είναι ίσο με τη διαφορά των εμβαδών των δύο κύκλων, δηλαδή $E_{\text{ΕΓ}} = \pi \cdot R^2 - \pi \cdot r^2$ ή $E_{\text{ΕΓ}} = \pi(R^2 - r^2)$.

        Άρα $\dfrac{E_{\text{ΕΓ}}}{E_{\text{ΑΕ}}} = \dfrac{\pi(R^2 - r^2)}{\pi \cdot r^2}$ ή $\dfrac{E_{\text{ΕΓ}}}{E_{\text{ΑΕ}}} = \dfrac{R^2 - r^2}{r^2}$.

        \item Με όμοιους συλλογισμούς προκύπτει ότι $E_{\text{ΑΔ}} = \pi \cdot \rho^2$ και $E_{\text{ΔΒ}} = \pi(r^2 - \rho^2)$.

        Άρα $\dfrac{E_{\text{ΔΒ}}}{E_{\text{ΑΔ}}} = \dfrac{r^2 - \rho^2}{\rho^2}$.
    \end{rlist}

    \item Από τα α)i) και α)ii), για να αποδείξουμε τη ζητούμενη ισότητα $\dfrac{E_{\text{ΕΓ}}}{E_{\text{ΑΕ}}} = \dfrac{E_{\text{ΔΒ}}}{E_{\text{ΑΔ}}}$ αρκεί να αποδείξουμε ότι $\dfrac{R^2 - r^2}{r^2} = \dfrac{r^2 - \rho^2}{\rho^2}$ ή $\dfrac{R^2}{r^2} - 1 = \dfrac{r^2}{\rho^2} - 1$ ή $\dfrac{R^2}{r^2} = \dfrac{r^2}{\rho^2}$ ή $\dfrac{R}{r} = \dfrac{r}{\rho}$.

    Η τελευταία ισότητα είναι αληθής, γιατί:

    Από εφαρμογή του θεωρήματος του Θαλή, το τρίγωνο ΑΔΕ που ορίζεται από τις ευθείες των πλευρών ΑΒ και ΑΓ του τριγώνου ΑΒΓ και την παράλληλη ΔΕ στην πλευρά του ΒΓ έχει πλευρές ανάλογες προς τις πλευρές του ΑΒΓ. Άρα $\dfrac{\text{ΑΔ}}{\text{ΑΒ}} = \dfrac{\text{ΑΕ}}{\text{ΑΓ}}$ ή $\dfrac{\rho}{r} = \dfrac{r}{R}$ ή $\dfrac{R}{r} = \dfrac{r}{\rho}$.
\end{alist}